% GENERATED FILE. Edit canonical lessons, then run npm run generate:books.
% canonical-source-hash: fnv1a64:34be6501c137be0e
% canonical-lessons: LA-C346-deus, LA-C346-divus, LA-C346-numen, LA-C346-mortalis, LA-C346-immortalis, LA-R346-review-339-346

\chapter{Deī et deae}
\label{ch:la-dei-et-deae}
\hlchaptermodality{346}

\begin{chapteropening}
\textbf{By the end of this chapter:} \emph{I can talk about the gods in stories: a god, a deified one, divine power, mortal and immortal.}

Complete the last lesson of chapter 346: I can talk about the gods in stories: a god, a deified one, divine power, mortal and immortal.
\end{chapteropening}

\section[deus, deī]{deus, deī — a god}

\label{lesson:LA-C346-deus}



\begin{quote}
Before the new one: say the Latin for the common people, then the Latin for a dictator.
\end{quote}

\subsection*{You'll want to know: deus, deī}
\textbf{deus, deī} — \textquotedblleft{}a god\textquotedblright{}.

\begin{quote}
Poētae dē deīs multa canunt.
\end{quote}

\begin{quote}
\emph{Poets sing many things about the gods.}
\end{quote}

Masculine, 2nd declension; plural \emph{deī} or \emph{dī}. Its feminine \emph{dea} is taught.

\begin{cousinweb}
\textbf{deus} is inherited from Old Latin \emph{deivos}, Proto-Indo-European \emph{*deywós}, \textquotedblleft{}god\textquotedblright{} (EtymDB). English \emph{deity} comes from its family.
\end{cousinweb}

\begin{grammarlens}[title={What you've built}]
One more word for telling, explaining and describing at B1.
\end{grammarlens}

\subsection*{Your turn}
\begin{itemize}
\raggedright
  \item \emph{Say it:} \emph{deus, deī}
  \item \emph{Say it:} \emph{deus, deī}, once more
  \item \emph{Say it:} a sentence of your own with \emph{deus, deī}
  \item \emph{Recall:} say the Latin for the common people, then the Latin for a dictator, then say \emph{deus, deī} again
\end{itemize}

\begin{tcolorbox}[breakable,colback=teal!4,colframe=teal!35!black,title={Before you move on}]
What does deus, deī mean? (\textquotedblleft{}A god\textquotedblright{}.) Say it once more.
\end{tcolorbox}
\section[dīvus, dīvī]{dīvus, dīvī — a deified emperor; a god}

\label{lesson:LA-C346-divus}



\begin{quote}
Before the new one: say the Latin for a dictator, then the Latin for a god.
\end{quote}

\subsection*{You'll want to know: dīvus, dīvī}
\textbf{dīvus, dīvī} — \textquotedblleft{}a deified emperor; a god\textquotedblright{}.

\begin{quote}
Post mortem imperātor dīvus factus est.
\end{quote}

\begin{quote}
\emph{After his death the emperor was declared a god.}
\end{quote}

Masculine, also an adjective, \emph{dīvus, -a, -um}, 'divine'. Coins call dead emperors \emph{dīvus}. English \emph{divine} is of its family.

\begin{cousinweb}
\textbf{dīvus} is usually traced to the same Old Latin \emph{deivos}, \textquotedblleft{}god\textquotedblright{}, as \emph{deus}; Lewis \& Short compare Greek \emph{dios}, \textquotedblleft{}divine\textquotedblright{}. English \emph{divine} comes from its family.
\end{cousinweb}

\begin{grammarlens}[title={What you've built}]
One more word for telling, explaining and describing at B1.
\end{grammarlens}

\subsection*{Your turn}
\begin{itemize}
\raggedright
  \item \emph{Say it:} \emph{dīvus, dīvī}
  \item \emph{Say it:} \emph{dīvus, dīvī}, once more
  \item \emph{Say it:} a sentence of your own with \emph{dīvus, dīvī}
  \item \emph{Recall:} say the Latin for a dictator, then the Latin for a god, then say \emph{dīvus, dīvī} again
\end{itemize}

\begin{tcolorbox}[breakable,colback=teal!4,colframe=teal!35!black,title={Before you move on}]
What does dīvus, dīvī mean? (\textquotedblleft{}A deified emperor; a god\textquotedblright{}.) Say it once more.
\end{tcolorbox}
\section[nūmen, nūminis]{nūmen, nūminis — divine power, divine will}

\label{lesson:LA-C346-numen}



\begin{quote}
Before the new one: say the Latin for a god, then the Latin for a deified emperor.
\end{quote}

\subsection*{You'll want to know: nūmen, nūminis}
\textbf{nūmen, nūminis} — \textquotedblleft{}divine power, divine will\textquotedblright{}.

\begin{quote}
Nūmen deae hunc lūcum servābat.
\end{quote}

\begin{quote}
\emph{The power of the goddess protected this grove.}
\end{quote}

Neuter, 3rd declension: the nod of a god, and so divine power. English \emph{numinous}.

\begin{cousinweb}
\textbf{nūmen} is built, says Lewis \& Short, on \emph{nuō}, \textquotedblleft{}to nod\textquotedblright{}, seen in \emph{innuō}: a nod, the nod of a god. English \emph{numinous} comes from it.
\end{cousinweb}

\begin{grammarlens}[title={What you've built}]
One more word for telling, explaining and describing at B1.
\end{grammarlens}

\subsection*{Your turn}
\begin{itemize}
\raggedright
  \item \emph{Say it:} \emph{nūmen, nūminis}
  \item \emph{Say it:} \emph{nūmen, nūminis}, once more
  \item \emph{Say it:} a sentence of your own with \emph{nūmen, nūminis}
  \item \emph{Recall:} say the Latin for a god, then the Latin for a deified emperor, then say \emph{nūmen, nūminis} again
\end{itemize}

\begin{tcolorbox}[breakable,colback=teal!4,colframe=teal!35!black,title={Before you move on}]
What does nūmen, nūminis mean? (\textquotedblleft{}Divine power, divine will\textquotedblright{}.) Say it once more.
\end{tcolorbox}
\section[mortālis, mortāle]{mortālis, mortāle — mortal}

\label{lesson:LA-C346-mortalis}



\begin{quote}
Before the new one: say the Latin for a deified emperor, then the Latin for divine power.
\end{quote}

\subsection*{You'll want to know: mortālis, mortāle}
\textbf{mortālis, mortāle} — \textquotedblleft{}mortal\textquotedblright{}.

\begin{quote}
Omnēs hominēs mortālēs sunt.
\end{quote}

\begin{quote}
\emph{All human beings are mortal.}
\end{quote}

Two endings, from \emph{mors}, 'death'. As a noun, \emph{mortālēs}: human beings.

\begin{cousinweb}
\textbf{mortālis} is built on \emph{mors}, \textquotedblleft{}death\textquotedblright{} (EtymDB): subject to death. English \emph{mortal} comes from it.
\end{cousinweb}

\begin{grammarlens}[title={What you've built}]
One more word for telling, explaining and describing at B1.
\end{grammarlens}

\subsection*{Your turn}
\begin{itemize}
\raggedright
  \item \emph{Say it:} \emph{mortālis, mortāle}
  \item \emph{Say it:} \emph{mortālis, mortāle}, once more
  \item \emph{Say it:} a sentence of your own with \emph{mortālis, mortāle}
  \item \emph{Recall:} say the Latin for a deified emperor, then the Latin for divine power, then say \emph{mortālis, mortāle} again
\end{itemize}

\begin{tcolorbox}[breakable,colback=teal!4,colframe=teal!35!black,title={Before you move on}]
What does mortālis, mortāle mean? (\textquotedblleft{}Mortal\textquotedblright{}.) Say it once more.
\end{tcolorbox}
\section[immortālis, immortāle]{immortālis, immortāle — immortal, undying}

\label{lesson:LA-C346-immortalis}



\begin{quote}
Before the new one: say the Latin for divine power, then the Latin for mortal.
\end{quote}

\subsection*{You'll want to know: immortālis, immortāle}
\textbf{immortālis, immortāle} — \textquotedblleft{}immortal, undying\textquotedblright{}.

\begin{quote}
Deī immortālēs numquam moriuntur.
\end{quote}

\begin{quote}
\emph{The immortal gods never die.}
\end{quote}

Two endings: \emph{in-} 'not' + \emph{mortālis}. \emph{Dī immortālēs!}: Good gods!

\begin{cousinweb}
\textbf{immortālis} joins \emph{in-}, \textquotedblleft{}not\textquotedblright{}, and \emph{mortālis}, \textquotedblleft{}mortal\textquotedblright{} (EtymDB). English \emph{immortal} comes from it.
\end{cousinweb}

\begin{grammarlens}[title={What you've built}]
One more word for telling, explaining and describing at B1.
\end{grammarlens}

\subsection*{Your turn}
\begin{itemize}
\raggedright
  \item \emph{Say it:} \emph{immortālis, immortāle}
  \item \emph{Say it:} \emph{immortālis, immortāle}, once more
  \item \emph{Say it:} a sentence of your own with \emph{immortālis, immortāle}
  \item \emph{Recall:} say the Latin for divine power, then the Latin for mortal, then say \emph{immortālis, immortāle} again
\end{itemize}

\begin{tcolorbox}[breakable,colback=teal!4,colframe=teal!35!black,title={Before you move on}]
What does immortālis, immortāle mean? (\textquotedblleft{}Immortal, undying\textquotedblright{}.) Say it once more.
\end{tcolorbox}
\section[(dialogue)]{Words from chapters 339–346 — a review}

\label{lesson:LA-R346-review-339-346}



\begin{quote}
Say the Latin for a battle and a battle line.
\end{quote}

\subsection*{Your turn}
Say these aloud:

\begin{itemize}
\raggedright
  \item a battle, a battle line, an army on the march, a raid, to fight it out
  \item weapons, a horseman, a foot soldier, a breastplate, a round shield
  \item to besiege, a siege, to attack, to take by storm, a battering ram
  \item to put to flight, to overcome, to cut down, to lay waste, a defeat
  \item a prisoner, to give up, a surrender, a treaty, a triumph
  \item ancestors, annals, a founder, an origin, great age
  \item a king, to reign, a tyrant, the common people, a dictator
  \item a god, a deified emperor, divine power, mortal, immortal
\end{itemize}

\begin{tcolorbox}[breakable,colback=teal!4,colframe=teal!35!black,title={Before you move on}]
A battle? (\textbf{proelium, proeliī}.) A battle line? (\textbf{aciēs, aciēī}.) An army on the march? (\textbf{agmen, agminis}.) A raid? (\textbf{incursiō, incursiōnis}.) To fight it out? (\textbf{dīmicō, dīmicāre}.) Weapons? (\textbf{arma, armōrum}.) A horseman? (\textbf{eques, equitis}.) A foot soldier? (\textbf{pedes, peditis}.) A breastplate? (\textbf{lōrīca, lōrīcae}.) A round shield? (\textbf{clipeus, clipeī}.) To besiege? (\textbf{obsideō, obsidēre}.) A siege? (\textbf{obsidiō, obsidiōnis}.) To attack? (\textbf{oppugnō, oppugnāre}.) To take by storm? (\textbf{expugnō, expugnāre}.) A battering ram? (\textbf{ariēs, arietis}.) To put to flight? (\textbf{fugō, fugāre}.) To overcome? (\textbf{superō, superāre}.) To cut down? (\textbf{caedō, caedere}.) To lay waste? (\textbf{vāstō, vāstāre}.) A defeat? (\textbf{clādēs, clādis}.) A prisoner? (\textbf{captīvus, captīvī}.) To give up? (\textbf{dēdō, dēdere}.) A surrender? (\textbf{dēditiō, dēditiōnis}.) A treaty? (\textbf{foedus, foederis}.) A triumph? (\textbf{triumphus, triumphī}.) Ancestors? (\textbf{maiōrēs, maiōrum}.) Annals? (\textbf{annālēs, annālium}.) A founder? (\textbf{conditor, conditōris}.) An origin? (\textbf{orīgō, orīginis}.) Great age? (\textbf{vetustās, vetustātis}.) A king? (\textbf{rēx, rēgis}.) To reign? (\textbf{rēgnō, rēgnāre}.) A tyrant? (\textbf{tyrannus, tyrannī}.) The common people? (\textbf{plēbs, plēbis}.) A dictator? (\textbf{dictātor, dictātōris}.) A god? (\textbf{deus, deī}.) A deified emperor? (\textbf{dīvus, dīvī}.) Divine power? (\textbf{nūmen, nūminis}.) Mortal? (\textbf{mortālis, mortāle}.) Immortal? (\textbf{immortālis, immortāle}.)
\end{tcolorbox}
